% TOP_PROBLEM: {"id": 20001039, "title": "A corrected colored reduction and a Sullivan-2 counterexample normal form", "category_id": 2, "category": {"id": 2, "name": "combinatorics", "display_name": "Combinatorics", "description": "Counting problems, graph theory, discrete structures.", "slug": "combinatorics", "order_index": 2, "created_at": "2026-07-31T15:26:25.671Z"}, "status": "partially_solved", "problem_number": "AIM-COMBINATORICS-0164", "set_id": 14}
% FIRST_POSTED: 2026-10-01
% ENTRY_KIND: statement_only
% UnsolvedMath ID 20001039; source catalogue code AIM-COMBINATORICS-0164
% !TeX program = lualatex
% Definition and sources only; this file is not an LLM solution attempt.
\documentclass[11pt,a4paper]{article}
\usepackage[margin=25mm]{geometry}
\usepackage{fontspec}
\setmainfont{lmroman10-regular.otf}[BoldFont=lmroman10-bold.otf,
 ItalicFont=lmroman10-italic.otf,BoldItalicFont=lmroman10-bolditalic.otf]
\usepackage{amsmath,amssymb,mathtools}
\usepackage{unicode-math}
\setmathfont{latinmodern-math.otf}
\AtBeginDocument{\let\setminus\smallsetminus}
\usepackage{xurl,hyperref}
\hypersetup{colorlinks=true,urlcolor=blue}
\setcounter{secnumdepth}{0}
\setlength{\parindent}{0pt}
\setlength{\parskip}{5pt}
\setlength{\emergencystretch}{3em}
\begin{document}
\section*{A corrected colored reduction and a Sullivan-2 counterexample normal form}\label{20001039}
{\small UnsolvedMath ID 20001039 · AIM-COMBINATORICS-0164 · Combinatorics · AIM Workshop Problem Lists}
\subsection{Definitions and mathematical statement}
Let $D=(V,E)$ be a nonempty finite oriented graph, meaning a simple
directed graph with neither loops nor opposite pairs of arcs. Define
\[
 N^+(v)=\{u:v\to u\},\qquad N^-(v)=\{u:u\to v\},
 \qquad d^\pm(v)=|N^\pm(v)|.
\]
Let $N_2^+(v)$ consist of vertices at directed distance exactly two
from $v$: they are endpoints of paths $v\to u\to w$, with $w\ne v$
and $w\notin N^+(v)$. The number of paths is irrelevant; only distinct
endpoints are counted.
The source's Conjecture 6.4 asks whether some vertex always satisfies
\[
                        |N_2^+(v)|+d^+(v)\ge2d^-(v).
\]
Thus the number of vertices reachable in one or two steps should,
for at least one starting point, be at least twice that point's
indegree. The two out-neighbourhoods are disjoint by definition,
so their sum counts distinct vertices rather than walks.

All graph sizes and all orientations of all underlying simple graphs
are included. In particular, the digraph need not be a tournament,
strongly connected, or balanced in its degrees. A vertex of indegree
zero satisfies the assertion automatically. The inequality concerns
one vertex chosen after the graph is given; an average inequality or
an estimate for two-edge walk counts would be a different claim.
The indegree on the right is retained exactly as in the workshop,
rather than replaced by the outdegree from Seymour's conjecture.
\subsection{Short English statement}
Must an oriented graph contain a vertex whose first and second out-neighbourhoods together are at least twice its indegree?
\subsection{Sources}
\begin{itemize}
\item[S1] ULAM.AI and the UnsolvedMath Contributors, supplied problems.json, record 20001039 (AIM-COMBINATORICS-0164), AIM Workshop Problem Lists. Catalogue identity and source transcription; CC BY 4.0.
\url{https://huggingface.co/datasets/ulamai/UnsolvedMath}.
\item[S2] American Institute of Mathematics, The Caccetta-Haggkvist conjecture, item 6.4. Original workshop question underlying AIM-COMBINATORICS-0164.
\url{https://aimath.org/WWN/caccetta/caccetta.pdf}.
\end{itemize}
\end{document}
